\documentclass{subfiles}
\begin{document}
    Among the many problems that algebraic geometry tries to solve,
        there is that of solving systems of homogeneous polynomials.
        That is, we'd like to use the tools provided by algebra and geometry 
        to study the zeros of polynomials. 
        Meaning we'd like to solve systems of the form 
        \[
            \begin{cases}
                f_{1} (x_{1}, x_{2}, \ldots, x_{m})  = 0 \\ 
                f_{2} (x_{1}, x_{2}, \ldots, x_{m})  = 0 \\
                \qquad \vdots \qquad \vdots \\ 
                f_{n} (x_{1}, x_{2}, \ldots, x_{m})  = 0 \\ 
            \end{cases}
        \]
        where all \(f_{i}, i = 1, 2, \ldots, n\) are either homogeneous polynomials
        over some projective space, or simply polynomials.

    We now introduce one of the fundamental results on the topic,
        which will be proven later on. 

    \begin{theorem}
        Let \(\F\) be a field algebrically closed. Then, a homogeneous system 
            of \(n\) equations admits non-trivial solutions if and only if 
            \(n \le m\), \(m\) being the number of variables in each polynomial.        
    \end{theorem}

    Another issue is that of describing the geometry of such solutions.
        In this context, the following classification is available.
        \begin{itemize}
            \item \(\F\) is an algebrically closed field: 
                \begin{itemize}
                    \item \textbf{dim 1:} the result \(\C* \subseteq \mathbb{A}^{2}\) is 
                        generalized. 
                    \item \textbf{dim 2:} several results by \emph{Castelnovo, Eriques, Severi, etc.}
                    \item \textbf{dim} \(\boldsymbol{\ge}\) \textbf{3:} many open questions.
                \end{itemize}

            \item \(\F\) is not an algebrically closed field:
                \begin{itemize}
                    \item Specific restrictions. For instance, the arithmetic algebric geometry 
                        used to prove Fermat's theorem (\(x^{n} + y^{n} = z^{n}\)).
                \end{itemize}
        \end{itemize}

    \subsection{Algebra reminders}
    \subfile{../sub/1.1 - Algebra reminders}

    \subsection{Noetherian rings}
    \subfile{../sub/1.2 - Noetherian rings}
    \clearpage
\end{document}
